\chapter{導論}
\section{文字與書寫系統}
首先，
我們需要對我們研究的對象——即文字與書寫系統——進行一個定義。
文字既然是語言的附庸，
我們很容易想到要從語言的角度出發來討論文字。
從歷時上來看，
文字是出於記錄資訊的目的出現的。
文字必然晚於語言而誕生，
我們當然期望意義和文字是同構的，
而理想狀況下我們已經有了意義和語言的同構，
那建立語言和文字的同構就成了一種方便的選擇。

對於一門具體的語言（英語、漢語等等），
文字就要建立一套具體的視覺符號和一門具體的語言之間的同構。
我們就可以理想地定義：與某門語言同構的一套視覺符號系統，
就是\Term{shu1xie3xi4tong3}{書寫系統}{writing system}。
當然，
現實沒有那麼理想。
於是，
我們有一個更弱的定義：能夠按照一定次序（往往是線性的）書寫某門語言的視覺符號，
就是書寫系統。

大家也許注意到了，
上面定義書寫系統時，
我總在強調「某門語言」。
這是因爲，
我想要把相互之間十分「相似」的書寫系統，
統稱爲\Term{wen2zi4}{文字}{script}。
這個詞就沒有什麼好的定義了，
但是通過例子可以直觀地理解：

\begin{ExampleBlock}
漢語：世界上有多種多樣的文化。

英語：There are many diverse cultures in the world.

德語：Auf der Welt gibt es eine große Vielfalt an Kulturen.

拉丁語：In mundo multae culturae diversae sunt.
\end{ExampleBlock}

我們看一下上面的句子。
英語所使用的那些字母顯然是符合書寫系統的定義的，
所以構成一種書寫系統，
我們不妨叫做「英文」；同理，
我們可以得到「德文」和「拉丁文」。
德文跟英文當然不一樣，
因爲他們所表達的語言都不一致。
但是，
我們顯然會自然地覺得，
這三種書寫系統是「相似的」，
他們的大多數符號好像都長得一樣，
儘管德文似乎多出了奇怪的字母：ß。
於是，
我們立即又說，
這三種書寫系統是同一種文字——我們統稱爲「拉丁文」。

我相信區分這兩個概念是有意義的，
儘管在漢語中我們都會使用「某某文」來稱呼之，
而且確實很多情況下不會區分得太嚴格。

以文字和書寫系統爲基本研究對象的學科，
就是\Term{wen2zi4xue2}{文字學}{graphemics}。
依照其研究的目的或着側重點，
可以區分出文字類型學、文字史、文字構型學、正字法等多種分支。

當然，
我們可以問出一個問題：文字真的就是語言的附庸嗎？首先這是一個過於激進的表述，
通常來說我們並不否認文字本身所存在的主動性。
這裏只不過是強調文字之作用不得不中介語言罷了。

如果只討論自然的文字，
我們往往會如此下論斷——我們於是也就這麼定義文字。
理論上來說，
文字只是與意義同構的一套符號，
沒有絕對的理由使之一定要經過語言。
我們只能說，
語言是文字自然產生比較方便的路徑。

歷史上來說，
語言的產生早於文字太多。
認知上來說，
人也是首先通過語言塑造了自己對意義的把握的——也就是說，
意義未必先於語言而存在。
這一點我們不作過多討論，
但是至少語言某種程度上反而能決定意義的構造。
所以，
當我們想要記錄已經被語言所塑造的意義時，
無論我們多麼想脫離語言，
都很難做到。

當然，
形式上的脫離並非實現不了。
譬如，
如果你選擇製造一種人工語言，
然後又特立獨行地使用表意文字作爲你的語言的書寫形式，
最可能的一種情況是，
這門「語言」可能只有你自己會使用，
你也許不會跟人說這門語言，
而更多會使用在你的筆記、日記等場合中。
這樣的話，
對這門語言來說，
作爲單純的視覺符號的文字就比作爲聲音符號的語言使用得更頻繁。
最後，
文字與其聲音形式脫鉤，
也就導致形式上這種文字脫離了語言。
但是，
上述的情況並不會是普遍的，
而且不會是本質的（其理由又涉及哲學討論了，
因此按下不表）。

還有一種更普遍的情形，
即標點、數字以及各種形式語言，
但顯然我們的論域已經排除了這些部分。

\section{文字類型學概論}
分類是重要的科學研究方法，
能夠將複雜而龐大的研究對象簡化爲可操作、可處理的不同類別，
有助於我們進一步展開研究。
因此，
對文字進行分類完全是有必要的。
分類很多時候是主觀的，
取決於我們關心什麼，
進而選擇什麼標準去進行分類。
對於文字來說，
自然且可操作的分類依據，
就是文字與語言之間所形成的映射關係。
於是我們狹義地定義：以文字和語言之間的映射關係爲依據研究文字的分類的學科，
就是\Term{wen2zi4lei4xing2xue2}{文字類型學}{typology of scripts}。
這個看法當然有很多問題，
但是可能是經典且十分好用的。
這一節對於文字類型學只是提綱挈領的介紹，
不會討論過多例外或者細節。

\subsection{語言系統的層級}

我會不斷強調文字作爲語言附庸的特性。
文字的功能是記錄意義，
而文字對意義的記錄必須中介語言，
討論這一中介的作用就是自然的了。
文字和語言的映射，
實際上是文字的基本單位與語言的單位的映射。
所以討論語言有哪些所謂的「單位」，
是十分有必要的。

語言作爲一種聲音符號，
當作爲聲音被發出時是具有物理性的。
語言作爲物理對象被研究時，
一般稱爲\Term{yu3yin1}{語音}{speech}。
人的發聲器官（請直覺地理解這個詞，
因爲過多的解剖學討論是偏離主題的）理論上能發出的所有聲音，
都可以稱作語音。
而研究語音的語言學分支，
稱爲\Term{yu3yin1xue2}{語音學}{phonetics}。
語音學研究的是你發出某個音時你的發聲器官是如何運作的，
以及你所發出的聲音具有怎麼樣的物理性質等等。

語音的最小單位是\Term{yin1su4}{音素}{phone}，
實際上可以理解成當你的發聲器官處於某種狀態或最短動程時所發出的聲音。
由於每個人的發聲器官不完全一樣，
而人也可以在一定限度內相對自由地調整自己發聲器官的狀態，
所以理論上音素總是無窮無盡的——或者說，
語音是一個連續統。
無窮無盡的音素是無法逐個研究的，
所以我們總是會嘗試限定一個範圍——比如只要你發音時兩個嘴脣必須接觸，
就叫做「雙脣音」，
而不去細究你到底閉得有多緊。
換句話說，
也就是先確立等價類，
再研究等價類的性質——所以音素實際上就是我們確立的最小必要的等價類。

但是，
當語言真正作爲語言，
即一個內部一致的符號系統時，
就不能夠只考慮其物理性質了。
我們更傾向於研究，
在一門語言中，
母語者是怎麼感受其所說語言的。
我們強調語言的結構屬性，
按結構主義語言學的看法，
一門語言不過是某個範圍內的人羣共同使用的聲音結構罷了，
而這一結構的特性完全僅僅取決於其中元素之間的關係。
這個關係往往體現在對立上，
比如說漢語中「蘋果」這個詞，
我們直接把文字對應在其語音形式上，
然後定義$\left[\text{APPLE}\right]:=\psi(\text{蘋果})$（即「蘋果」所對應的意義）。
實際上，
假定我們說服了所有漢語母語者不再使用「蘋果」這個詞來表示[APPLE]，
而使用「甜紅果」來表示，
那麼「甜紅果」就成爲了[APPLE]在漢語中所對應的詞，
這完全是可以的。
但是，
如果讓所有人都把[APPLE]都叫「香蕉」，
聽起來就沒那麼能接受了。
因爲我們的語言中已經有了「香蕉」這個詞，
用來表示[BANANA]。
如果我們再把[APPLE]也叫「香蕉」，
就完全亂套了——因爲這兩種事物都是水果，
而且都是很常見的水果，
也很少有人認爲兩種水果差不多（當然這個討論並不絕對，
當作一個淺顯但不準確的例子就好）。
通過上面這個例子我是想說明，
用什麼來表示[APPLE]並不重要，
只要我們覺得有必要與[APPLE]區分開來的東西不會用到同一個詞就行了——也就是說，
一個詞是否合法，
完全取決於他能夠與那些他需要區別開來的事物區別開。

在語言中，
如果兩種元素被母語者覺得應該區分開了，
無論是出於什麼目的，
我們都傾向於分爲兩個元素，
而不是歸併爲一種；反之，
如果母語者覺得區分不開，
那也就沒必要認爲是兩種了。
比如漢語中沒有濁的[b]，
如果你把漢語拼音的b發成[b]，
完全沒問題——只要有跟漢語拼音p足夠的區分度就好。
實際情況下，
我們也能夠發現漢語的所謂清塞音確實比法語之類的清塞音更濁一點，
大概就是漢語沒有濁塞音的緣故。
上面說了那麼多，
我無非只是想要強調結構主義語言學的一個基本主張：\textit{語言作爲一個共時的系統，
其單位的價值完全取決於該單位在系統之間的差異與對立。
}

這樣，
語言非物理的部分，
其最小的單位就是\Term{yin1wei4}{音位}{phoneme}。
音位可以理解成音素在語言中的等價類，
在共時的系統內，
等價類中的不同音素視作一個音位——只要一門語言中不同的音素沒有區分的必要，
就應當視爲一個音位，
這些不同音素稱爲\Term{yin1wei4bian4ti3}{音位變體}{allophone}。
由音位所組成的系統稱爲\Term{yin1xi4}{音系}{phonology}，
音系就是一門語言中可能出現的所有音位的集合。
利用數學語言，
對於一門語言中的語音$P$，
定義其上的等價關係$\sim$，
則音系定義爲商集$P/\sim$，
任意$\bar{p}\in P/\sim$即一個音位，
顯然$P=\coprod\bar{p}$。
從上面的討論可以看出，
之所以說音位不是完全物理性的，
就是因爲這一等價關係某種程度上是人爲規定的。
研究音位，
或者等價地也是研究音系的學科就是\Term{yin1xi4xue2}{音系學}{phonology}。
音系學研究的是一門語言的音系是怎樣的，
即有哪些音位等等。

需要注意的是，
語言作爲聲音所具備的音色、音高、音強、音長等特徵，
都可以作爲區分音位的依據。
如果我們僅僅就音色來討論音位，
實際上僅僅涉及了音段。
\Term{yin1duan4}{音段}{segment}實際上指的就是以音色爲主要區分依據的音位，
以其他物理性質來區分音位的則是\Term{yun4lv4}{韻律}{prosody}。
韻律就包括聲調、重音、節奏、響度等等，
這些特徵往往不會僅僅出現在一個音段中，
而是常常表現在連續的幾個音段上，
因此也常常稱爲\Term{chao1yin1duan4}{超音段}{suprasegmental}。
不過，通常來講，「音位」這個詞總是傾向於指音段音位——最多再包括重音和聲調。

音段包括元音與輔音。
一般來說，
\Term{yuan2yin1}{元音}{vowel}指發音時氣流能自由從口腔通過不受阻礙的音；反之，
如果發音時在口腔的任何一個部位受到阻礙，
就是\Term{fu3yin1}{輔音}{consonant}。
對韻律的詳細討論對我們來說是不必要的，
我們只提一下裏面比較重要的幾個概念。

粗淺地來說，
在語音中比其他部分在聽覺上更爲顯著的某些部分即是\Term{zhong4yin1}{重音}{accent}。
這種意義上的重音主要體現在音強更強，
也稱爲\Term{zhong4du2xing2zhong4yin1}{重讀型重音}{stress accent}。
還有另一種特徵我們也稱爲重音，
即當一門語言中不同部分有音高的對比，
但音高的對比十分簡單時，
這種音高對立我們稱爲\Term{yin1gao1xing2zhong4yin1}{音高型重音}{pitch accent}。
其中，
音高對比比較簡單指的是，
往往最多只有高、中、低三個音高的對比，
沒有相當分明的上升、下降或曲折（當然自然的升降仍然會有）。
日語、吠陀梵語都被認爲具有音高型重音。

如果一門語言有相對複雜的音高變化，
即具有上升、下降和曲折，
以及多種不同音高的對比，
而這種音高變化具有比較明確的特定種類，
並承載區別意義的功能時，
我們即稱之爲\Term{sheng1diao4}{聲調}{tone}。
漢語具有聲調，
因此大家對此已經相當熟悉了。

\Term{yin1chang2}{音長}{length}即發音的長短。
如果音長並非系統性區分意義的因素，
我們不會將之視爲區分音位。
如果一門語言中有明顯的區分意義的長、短兩個相同音色的音素，
我們才會將之視爲音位的差別。
例如，
在漢語中拉長某個字除了表達強調或者某種情感之外，
不會被我們認爲改變了什麼意義。
反而，
在德語中，
Stadt和Staat音長的不同，
則是完全區別意義的。
不過實際操作時，
我們一般把長短元音歸爲元音的區別，
把長短輔音歸到兩個音節中，
而不會單把音長作爲非音段拿出來說。

音位可以按某種規則組合成\Term{yin1jie2}{音節}{syllable}，
音節就是人能夠自然感覺出的最小語音結構單位。
實際上，
音節是很複雜的單位，
從我們模糊的定義中就能夠隱約感覺到。
我們也不對音節作過細的討論。
音節一般要有一個\Term{he2}{核}{nucleus}，
這是音節中最顯著的部分，
一般由元音充當，
少部分輔音也偶爾可以充當。
除了核之外，
音節的前前後後還可以加上許多音段作爲次要成分，
也可以具有超音段成分。
我們直接看一個例子來避免複雜的討論：漢語的「牀」，
漢語拼音是chuáng，
其中音節核是a，
而ch和u是前面的次要成分，
ng是結尾的次要成分，
聲調則是超音段成分。

上面的音位和音節，
都是未必具有意義的。
把一個「牀」拆出一個音位ch，
你當然不能說ch有什麼具體的含義；philosophy拆出一個lo，
那確實很難說這個lo有什麼意思。
我們只介紹態位和詞兩個能具有意義的語言單位。

具有詞彙意義或語法意義的最小語言單位就是\Term{tai4wei4}{態位}{morpheme}。
對於eater這樣的詞，
eat和-er是兩個態位。
當然，
和我們的「音位」類似，
態位同樣是某種等價類特性的概念，
其下位具體的不同形式我們稱爲\Term{tai4su4}{態素}{morph}，
例如意大利語中的allargamento和approfondimento，
al-和ap-是同一個態位（注意，
一定要使用語音形式；但幸而意大利語的語音形式是和書寫系統比較對齊的），
因爲無論是語義還是歷史來源，
這二者都是一致的。
但是，
這兩個不同的形式是兩個態素。
類似的，
我們也稱同一個態位所轄的不同態素爲其\Term{tai4wei4bian4ti3}{態位變體}{allomorph}，
上面的al-和ap-就是同一態位的不同態位變體。
研究這些內容的語言學我們稱爲\Term{xing2tai4xue2}{形態學}{morphology}。

我們仍然引入\Term{ci2}{詞}{word}的概念，
但最好從直觀上去理解，
因爲這個概念多少有些模糊——不同語言對詞的理解似乎都有一點不同，
導致我們只能大概給出一個定義：即詞是能獨立運用的最小語言單位。
我們不對此作太多討論，
只需要記住，
言及態位和詞時，
我們都意指帶有意義的語言單位。

我們一共介紹了音位、音節、態位三個語言單位，
而語言的運作機理，
實際上是通過音位或音節的組合，
定義一個到意義的映射，
從而使這個組合成爲態位，
再按照意義的組合，
定義態位的組合方式，
進而形成更大的單位來表達複雜意義。



\subsection{文字的類型}
對於一種文字，
我們或許可以自然地找到其基本單位：比如漢字中的一個個「字」、拉丁文中的一個個「字母」。
但是事實沒有那麼容易。
我們在下面舉出一些例子：

\begin{ExampleBlock}
漢字：「河」、「江」爲什麼只算成兩個基本單位？爲什麼不能拆成「氵」、「可」、「工」三個基本單位？

拉丁文：qu基本完全不拆開，
爲什麼不能算成一個基本單位？a和á應該算一個還是兩個基本單位？

阿拉伯文：一個字母，
如\AR{ت}，
在詞首寫成\AR{تـ}，
在詞中寫成\AR{ـتـ}，
在詞末寫成\AR{ـت}，
爲什麼不能視爲四個基本單位？

天城文：字母\SA{क}和\SA{का}，
應該視爲兩個不同的基本單位，
還是應該拆爲\SA{क}和\SA{ा}？

…………
\end{ExampleBlock}

很可惜，
我們幾乎無法簡單地回答上面的問題，
而只能首先從直覺，
或者使用者的感覺上去處理。
比如從漢字使用者的角度來說，
把「字」視爲基本單位是自然的，
我們確實可以說出很多的理由：比如部件多半不能單獨使用、不同的部件組成字的規律並不是普遍的、你也不能夠隨意組合部件等等。
但是這些理由顯然並不普適，
於是我們最好還是直接接受現成的結論，
將「基本單位」的認識統一爲某種習慣上的約定。
實際上，
當我們接受了這個習慣上的約定，
再對類型學有更深入的瞭解之後，
我們可能會有一些新的看法。


不管怎樣，
我們首先定義文字中最小的區分性單位爲\Term{zi4wei4}{字位}{grapheme}。
這個字位在文字中的地位和音位在語言中的地位相當相似，
譬如對於一個字母a，
我往往不在乎到底是寫得粗（\textbf{a}）、寫得細，
正體還是斜體（\textit{a}），
是單層（ɑ）的還是雙層（a）的——只要在文字這個系統中，
能夠與其他的字位形成對立，
就足以把其自身確立爲一個字位了。
按我們習慣上的理解，
漢字中的「字」、拉丁文中的「字母」等等就是字位。
我們通常把文字中各種不同寫法稱爲\Term{zi4su4}{字素}{graph}，
或者排印學中更常稱爲\Term{zi4xing2}{字形}{glyph}，
而屬於同一個字位的不同字形則是\Term{zi4wei4bian4ti3}{字位變體}{allograph}。
譬如說，
a、ɑ、\textit{a}、b、\textit{b}都是不同的字形，
而字形a、ɑ、\textit{a}是同一個字位$\langle\text{a}\rangle$的不同字位變體，
b、\textit{b}則是字位$\langle \text{b}\rangle$的不同字位變體。
平時我們所說的「書體」、「字體」，
實際上都是字形的差異。
研究不同字形的歷史演變一般是文字史的研究內容，
但現代字形往往是排印學的主要研究對象——這已經不是單純的文字學了。

我們很容易觀察到我們所定義的概念在邏輯上的對應關係：
\begin{ExampleBlock}
音系：音素——音位——音位變體

形態：態素——態位——態位變體

文字：字素——字位——字位變體
\end{ExampleBlock}

定義了上述概念後，
我們就可以研究文字與語言之間映射的關係了。
語言的各個單位——音位（元音、輔音、超音段）、音節、態位，
都可以成爲字位的像。
在上述的單位中，
映射到音位還是音節，
都會導致獨立的字位往往不表達意義，
我們就稱爲\Term{biao3yin1wen2zi4}{表音文字}{phonogram}。
而如果存在相當一部分字位映射到態位或詞，
則同時具有了表達意義的能力，
那我們就稱之爲\Term{biao3yi1wen2zi4}{表意文字}{ideogram}。

最容易想到的是，
如果字位能夠和音位一一映射，
那必然是相當方便的。
現在世界上絕大多數在用的文字都是這一種邏輯，
因爲一方面這樣的文字字位少，
易學易記易用；另一方面，
這個方法是最爲普適的，
任何語言都可以這樣去創制文字，
或者改造別人已有的文字爲己用。
這種字位與音位一一對應的文字我們就稱爲\Term{quan2yin1wei4wen2zi4}{全音位文字}{alphabet}。

世界上有一些語言的元音並不發達，
或者往往只承擔語法意義。
使用這樣語言的人往往會重視輔音多於元音，
而創製出一種輔音地位遠遠高於元音的文字。
這樣的文字往往會擁有映射到全部輔音的不同字位，
但卻省略掉很多元音所對應的字位。
有意思的是，
世界上最早的表音文字——腓尼基文——就是沒有元音字位的。
我們把這樣的文字稱爲\Term{fu3yin1yin1wei4wen2zi4}{輔音音位文字}{abjad}。
現代常見的文字中，
阿拉伯文和希伯來文都屬於輔音音位文字。
這兩種文字都可以省略大多元音字位，
只保留了一部分不可省略；如果一定要標出元音的話，
可以另外用特定的符號標示出來。

有的語言中，
有某個元音（一般是a）尤其常見，
而其他元音的數目則相對來說比較少。
使用這樣語言的人可能就會在大多數情況下省略這個元音所對應的字位，
而保留其他元音以及所有輔音所對應的字位。
這樣的文字就是\Term{yuan2yin1fu4biao1wen2zi4}{元音附標文字}{abugida}。
世界上現存的大多數元音附標文字都有共同的起源——婆羅米文。
由婆羅米文衍生出來的一大批文字，
我們都稱爲婆羅米系文字。
天城文就是一種常見的元音附標文字，
這種文字中對應的a往往不會寫出來，
但其他的元音都有特定的符號來表示（除了元音前沒有輔音的時候）。
藏文是一種更徹底的元音附標文字，
所有的a都不會單獨寫出來，
而其他所有元音則都有相應的符號表示。

對於一些音節結構比較簡單、音節數目不太多的語言，
直接把字位映射到音節上是一種很好的想法。
我們把直接把字位映射到音節的文字稱爲\Term{yin1jie2wen2zi4}{音節文字}{syllabary}。
假名就是典型的音節文字，
其所對應的日語開音節相當多、音節結構相對簡單，
很適合使用音節文字來表示。

上面所述的四種文字即我們所稱的表音文字。
所描述的語言特徵與文字特徵的對應關係往往不是絕對的，
只是一種簡化的解釋模型。
譬如，
我們都知道，
雖然假名足以表示日語，
但總會有一些其他的因素導致日語沒有辦法只使用假名。
而像是諾蘇彝文，
則是經過人爲規範化並進行刪改的結果，
其成因不是自然的。
輔音音位文字也可以通過增補元音符號成爲全音位文字，
例如維吾爾文。
另外，
一種語言採取某種文字，
更多時候是出於歷史原因而不是自然的原因。
更多具體的例外和案例，
我們會在以後的章節進行討論。

上面對表音文字的討論，
缺陷也是很明顯的。
我們的文字類型學，
往往冠名爲對文字的討論，
可是與語言的對應又往往落到書寫系統上。
這必然會招致一定的矛盾，
因爲文字和書寫系統不必然擁有相同的性質：一種文字有時在不同書寫系統中可以對應到不同的語言層級（比如阿拉伯文與維吾爾文），
那就不應該一概而論；而同一套書寫系統可能不只包含一種文字（比如日文使用了假名和漢字）——雖然這種情況很好解決，
只要分開討論就可以了。

另一方面，
儘管在我們的理想模型下字位與語言單位的對應應該至少是一個單射，
但實際情況中這個映射可能甚至不良定，
就更不必嘗試去討論是否單、是否滿了。
拿英文來舉例，
同樣的字位a完全可能對應兩個不同的音位。
有人可能說這可能是由某種拼寫環境條件決定的，
比如cat和cate中a的發音不一致實際上是由於後者末尾的e；也有人會說可能得考察歷史來源，
比如have和gave儘管拼寫類似，
但來源有異，
因此a的讀音不同——前者尚且還是能說通的，
我們只要納入拼寫環境條件作爲限制因素就可以了；後者則不應當作爲解釋，
因爲共時系統的規則不必求諸歷時，
歷史來源不是共時系統內部可用的規則性限制（用漢字構型來論述其類型也是犯了同樣的錯誤）。
因此，
從我們看見的結果來說，
由於歷史原因導致的映射不良確實會破壞我們的理想，
但這應該留待正字法學去討論，
文字類型學——至少是我們比較保守框架下的類型學——應該暫時選擇性忽略這些。

除了上述的類型，
我們顯然還能夠看到一些雖然並不主流，
但無法劃入上面框架內的表音文字：例如有些文字，
其字位映射到音位以上，
但還不到音節的程度，
我們稱這樣的文字爲\Term{ban4yin1jie2wen2zi4}{半音節文字}{semi-syllabary}。
注音符號就是一種典型的半音節文字（有人認爲注音符號不屬於文字，
但其確實符合文字的定義），
其將音節結構劃分爲聲、介（韻頭）、韻（韻腹和韻尾）、調四個部分，
字位則分別對應這四個部分裏面的單位——即一個字位可能代表一個音位，
例如「\ZY{ㄎ}」表示漢語拼音k；也可能表示一個音位組合，
例如「\ZY{ㄤ}」表示漢語拼音ang，
實際上是兩個音位。
無論如何，我們將映射到音節以下單位的，包括輔音音位文字、元音附標文字、全音位文字在內的文字類型，
也統稱爲\Term{yin1duan4wen2zi4}{音段文字}{segmental}。

現在我們來討論表意文字。
實際上，
現在所能見到的成熟的表意文字只有一種，
即\Term{yu3su4wen2zi4}{語素文字}{logogram}，
字位基本對應到態位層面。
但是，
實際上很多所謂的表意文字，
總是會形成一個混合系統，
即有部分符號表音，
而部分符號表意——我們仍然把這種文字視爲語素文字，
否則會使我們的討論過於細節、過於複雜。
那麼，
修正後的語素文字的定義就是，
存在相當一部分字位直接映射到態位（或詞）層面的文字。

實際上，
「語素文字」這個詞，
已經足夠把所有可能的表意文字囊括進去了。
雖然「表意文字」自稱「表意」，
但文字不能直接越過語言的中介表意，
不應當誤解「表意」這個名詞。
後文中，
我們直接把語素文字和表意文字兩個概念等同起來。


與「表意文字」相關的，
還有在中國學界很流行的兩個名詞，
一是「象形文字」，
一是「意音文字」。
這兩個名詞實際上是同一個問題，
即用歷時分析代替共時分析，
或者混淆了構型學和類型學。
所謂「象形文字」，
其依據不過是在文字產生初期往往會去描摹事物的狀貌，
但是首先純通過「象形」是不可能表示各種抽象概念與名詞之外的詞的，
也就不能夠成爲文字中的主導；其次，
「象形」最多只能是文字構型的方法，
即描摹事物來表達相應的概念，
是從文字歷時性的構成來進行的分析，
不能代替共時分析。
「意音文字」的主張者無非認爲在漢字內部也有「聲符」，
但是這也是歷時的漢字構型，
並不是說我可以隨意地在共時系統中選擇純表音的構件（沒錯，
「聲符」甚至不是字位）——所以也不能夠作爲類型學的分類依據。


我們再對上述的概念進行總結（我們把被省略的部分稱爲地位低）：
\begin{ExampleBlock}
全音位文字：字位一一映射到音位

輔音音位文字：字位映射到音位，
但輔音地位高於元音

元音附標文字：字位映射到音位，
但某個元音地位低於其他元音

音節文字：字位映射到音節

語素文字：字位映射到態位或詞
\end{ExampleBlock}

\section{文字發展概論}
對文字的類型學有一定瞭解之後，
我們有必要來討論一下文字史。
顧名思義，
研究文字的歷史演變的學科，
就是\Term{wen2zi4shi3}{文字史}{history of writing systems}。但是，
這個描述顯然不夠清楚。
與其他所有的歷史一樣，
文字史也不會是簡單的、線性的發展歷程，
而是有着許多不同的層次。我們在這一節，
不可能覆蓋所有的文字歷史，
而更希望以一個宏觀的、理論的角度
去探索文字演進的歷程。
而我們在這裏討論的，實際上只是文字史中的一部分，即文字在不同族羣間流變、演化的歷史。從這個角度來看，我們將文字發展類型分爲三類：
獨立型、他源型和競爭型。
在獨立型中，我們也許會簡單討論一個文字內部自然沿時序變化的過程，但這不是我們的重點。

\subsection{獨立型}
首先是最簡單的情況。
對於文字史來說，
最容易想到的是一門語言的使用者獨立、自然地發明了一種文字，
而這種文字就線性獨立發展，
不受到任何外來因素的影響。
我們不妨稱之爲\Term{du2li4xing2}{獨立型}{independent type}。

一般來說，
自然的文字發明不會是一瞬間的，
而往往經歷了一個很漫長的過程。
一開始，
人們可能只會想着用簡單的圖畫去描摹事物的樣貌，
而並不想着去
記錄語言。
這往往比較像是「助記」的符號，
協助他們去回憶起
重要的資訊，
而不一定需要完整記錄語言，
也就當然不能夠稱爲文字。
我們稱這種通過描摹樣貌等方式來記錄特定事物的視覺符號爲\Term{tu2fu2}{圖符}{pictogram}。
圖符在某種意義上接近我們傳統理解的「象形文字」，
但圖符也不僅僅通過象形來表示。
廣義上來看，各種巖畫、刻符等，也可以算作圖符。

當社會生活進一步複雜化，簡單的圖符就開始向文字轉化。
這應當是一個很漫長的過程，很多發展到後期的圖符甚至能夠基本確定地對應一段語言。
例如加拿大奧吉布瓦人（Ojibwe）曾使用一種相對複雜的圖符組合來表達意義，
這些圖符組合使用簡單的圖騰和線條，
往往能夠被解讀成差不多的意思，
在某種程度上已經具備了文字的一些特徵。
但這些圖符仍然不能按語詞的順序去書寫語言，
其對應往往相對比較寬泛，
只能從整體出發來理解，
因此並不能夠視爲文字。

但是，當圖符逐漸開始與語言掛鉤，形成一個個單位來表達語言中的相應成分時，我們就說文字產生了。
文字的產生似乎還有很多謎團，因爲我們似乎沒有真正發現過完整的由圖符到文字的演進序列。
不過，文字既然不可能憑空產生，那推斷其由圖符轉換而來則是有道理的。

早期，成熟的文字往往不會單純表音。我們所能夠發現的最早的文字，基本上都是表意文字系統，而且越早期的文字，就越帶有圖符的特徵，象形程度也就較高。
目前所知的最早的文字，應當是公元前四千紀左右產生的蘇美爾楔形文字和埃及聖書文，
都屬於表意文字，在產生初期都具有相當明顯的象形特徵。

文字產生以後，如果一直沒有中斷的話，所經歷的應當是不斷的抽象化的進程。
也許有人會認爲，簡化可以作爲文字發展的一大趨勢。
但這個詞多少有些容易讓人誤會，因爲很多時候，文字的發展是總體上的簡化和部分的複雜化長期共存的。
而且複雜化的部分，也容易大到不可忽視的程度。
於是，我們更傾向將文字的發展描述成適應其書寫環境和文化需要下的抽象化。

書寫環境對文字的影響是相當顯著的。正如楔形文字由早期的圖符化特徵明顯，到演變成我們熟悉的「楔子」形，
就是適應其蘆葦杆與泥板書寫的需要。
而商代產生的漢字，
在逐漸的演變中也適應書寫而變得筆畫化、平直化。

文化需要則是一部分人爲因素對文字發展的影響，且這一影響往往是較爲守舊的。
我們能夠在漢字中找到許多例子，譬如比較正式的青銅器上，其字形可能比甲骨文還更象形。
而秦始皇的書同文，也沒有選擇新興的隸書或其他書體，反而選擇了保守的秦篆。
當然，我們也不排除更爲「激進」的文化需要對文字的影響。
例如，漢字的今草就是一種相當人爲形成的書體，
其早就超過了日常書寫的範疇，而成爲一種對藝術表達的追求，
也應當屬於文化需要影響文字發展的一部分。
但是，這種影響似乎往往很難成爲主流，因爲比自然演進更激進的人爲影響往往會造成文字的可識讀性變差，
實際上是不利於其推廣的。

抽象化則反而應當是文字作爲符號而發展的更本質的特徵。
符號與意義的映射，本身就應當有着隨機性。
所謂早期文字的象形成分，不過是某種便於發明符號而採取的權宜之計。
隨着社會生活的逐漸複雜化，文字符號所需要表達的內涵也不斷抽象化。
那麼，如果文字不往抽象化的方向演進，就必然不能夠滿足社會生活的需要。
因此，文字就會自然地往抽象化的方向發展，與其所涵攝的意義的關聯也就越來越難以發現了。

我們很難斷定獨立型文字在後期的發展趨勢，畢竟完全符合獨立型定義的文字，在歷史上是幾乎不存在的——或者說往往只有自源文字的早期才會符合。
這不難理解，因爲人類歷史是交流的歷史，不太可能存在長期發展而又完全獨立的文化。
即便是美洲文明，也終究會被外來文明所影響——但美洲本身的社會發展水平一直不高，文字要很晚才出現，還沒有充足的發展時間就遭受了破壞。

不過，漢字和雲南彝文給出了近似後期發展的獨立型文字的案例。雖然這兩門文字也並非與外界隔絕，但其發展相當穩定、連續，某種程度上類似於獨立型文字。
這兩種文字最終定型爲語素文字，沒有再進一步發展。
當然，彝文的情況比較特殊。標準的涼山彝文在經過人爲的整理之後，就變成一種音節文字了。
不過雲南彝文則沒有經過這樣的整理，也就保留了語素文字的性質。

至於另外的一些歷史悠久的文字，要麼是被外來文字取代，要麼是在不斷的傳播中變成了表音文字。
埃及聖書文發展到世俗體時，迎來了托勒密王朝。希臘語對埃及語的取代，也就宣告了希臘文這一可能是聖書文後裔的文字，最終取代了聖書文的直繫後裔世俗體。
當然，也有一部分世俗體被吸收到一個新的體系中，這個體系的主體來自希臘文，在後來成爲了唯一埃及語現存後裔科普特語的文字科普特文。
楔形文字幾經轉手，在阿契美尼德王朝時期就已經成爲了一種音節文字，最後直接被亞蘭文取代。

如果一定要總結出一條線性的發展道路，我們會說語素文字有發展爲音節文字的傾向，而音節文字則容易發展出音段文字。在音段文字中，全音位文字和元音附標文字又相對
穩定，而輔音音位文字是有發展爲全音位文字的趨勢的——
但說實話，我很不希望有這樣線性的總結，因爲這會多多少少掩蓋掉文字發展中實際存在的複雜性。
我們這一節的重點，也不是探索一種文字是怎麼內部迭代的。

\subsection{他源型}
可以說，世界上絕大多數文字都是來源於更早的文字。他們或者直接拿別人的文字稍加刪改即爲己用，或者在別人文字的基礎上或啓發下創製出一種全新的文字。
這些來自於更早的其他文字的文字，我們稱爲\Term{ta1yuan2xing2}{他源型}{externally derived type}。與他源型相對，自我起源的文字就是\Term{zi4yuan2xing2}{自源型}{primary type}。
落實到討論文字史類型時，需要強調他源型是自己本身沒有文字，或者完全放棄了原有的文字，然後才在別人文字的基礎上有了自己的文字；
自源型顯然包括獨立型，但獨立型另外涵攝了幾乎不受外來文字影響的含義。

在他源型文字中，如果一個族羣直接將其他族羣的文字拿來用，而更改很少，可以視作同一種文字時，我們稱之爲\Term{shi4ying4xing2}{適應型}{adapted type}。
適應型的情況在歷史上相當常見。
在借用時，適應型連相應字母對應的發音都沒有太多改變。
例如，英語放棄了原有的如尼文，而改用拉丁文，實際上就屬於適應型。英語幾乎是把拉丁文原封不動拿來用，只增補了少量字母。
這種情況下，我們不會將之視爲兩種文字，即適應型不產生新文字。
那麼，在考察文字史時，這種發展類型常常可以忽略——也許文字史應該更關心產生新文字的進程，而文字系統內部的變化，似乎比較適合拿到下一節來講。

另外還有一種他源型文字，要麼只是借用其他文字的形，而形成一套完全不同的系統；要麼只是受到了其他文字的啓發，從而獨立創造出一種文字。
這種進程是會產生新的文字的，因爲文字的邏輯和語文映射關係已經完全被顛覆了。我們稱這種文字發展爲\Term{yan3sheng1xing2}{衍生型}{derived type}。
出於上面已經說過的考慮，我們也常常把他源型和衍生型等同起來。
衍生型在歷史上也不少見，例如羅馬人借用希臘文發明了拉丁文，女真人借用漢字發明了女真文，藏族人借用天城文發明了藏文，等等。

他源型是文字史上相當重要的一種文字發展類型，也是最普遍的一種。研究他源型文字的發展，需要我們結合歷史學的方法，將文字與文字之間比較後排出演進序列，這樣之後
我們就能夠找到文字是如何逐漸發展，如何從一點逐漸擴散的了。

歷史上，往往有某種文字，其對世界的其他文字產生了相當深遠的影響，有許多文字都是由這種文字衍生而來，我們就會把這所有的文字都歸爲一個系列。

例如，漢字在整個東亞有着舉足輕重的地位，
不僅啓發了西夏文、女真文、契丹文、假名、諺文等等一系列文字的誕生，也作爲適應型應用到越南語、壯語、日語、朝鮮語等語言中。
我們就可以將由漢字衍生而來的所有文字，連同漢字（包括其在不同語種中的變體）一起，稱作\Term{han4zi4xi4wen2zi4}{漢字系文字}{Sinographic scripts}。
當然，漢字的影響是與中國的影響不可分割的。除了文字，中國也在文化、語言等各個層面影響着東亞、東南亞的廣大區域，
這些區域我們統稱爲\Term{han4zi4wen2hua4quan1}{漢字文化圈}{Sinosphere}。

而印度在南亞、東南亞的影響，也和漢字在東亞的影響一樣不可低估。婆羅米文可能衍生自亞蘭文，產生於吠陀時代晚期到列國時代早期的印度。
經過漫長的發展，婆羅米文衍生出了天城文、藏文、孟加拉文等諸多不同的文字，這些文字我們統稱爲\Term{po2luo2mi3xi4wen2zi4}{婆羅米系文字}{Brahmic scripts}。
類似地，受印度文化影響的廣大區域，也可以統稱爲\Term{yin4du4wen2huan4quan1}{印度文化圈}{Indosphere}。

\subsection{競爭型}
還有一種比較特殊的文字發展類型，即一種文字的使用者受到外來文字的影響，但又不拋棄其原本的文字；或者原本使用其他族羣文字的族羣，又發明了自己的文字，
這兩種文字都沒有完全取代對方。這種文字發展類型我們稱之爲\Term{jing4zheng1xing2}{競爭型}{competitive type}，其結果要麼是融合，要麼並存，要麼形成混合系統。
競爭型文字發展的邏輯在於，有兩種或兩種以上的文字系統同時存在於同一族羣中，而又沒有相互完全取代。

競爭可能帶來融合，形成不同於競爭雙方的新文字。這種情況可能與衍生型文字十分類似。例如，在托勒密王朝時期的埃及，存在世俗體與希臘文兩種文字的競爭，
雖然希臘文佔上風，但在局部競爭的條件下，融合成了一種新的文字，即科普特文。

競爭也可能帶來並存，競爭的兩套文字都可以使用，但不會混合使用。例如，在現在的蒙古國，蒙古文與西里爾文即處於並存狀態。在經過長期的西里爾化之後，
蒙古國民衆已經習慣了西里爾蒙古文。但蒙古國政府近年來又在嘗試恢復傳統蒙古文，也就使得蒙古國現在實際上並行兩套文字。
當然，我們也知道，兩種文字的並存只是短暫現象，歷史的發展終將作出選擇。

競爭也可能會帶來混合系統。比較典型的實際上是日本，其在過去只使用漢字，在假名產生後，則使用漢字和假名共同書寫的混合系統。
儘管日本當局多次嘗試限制甚至棄用漢字，但其音節結構簡單、漢語借詞豐富，導致棄用漢字必然會帶來各種各樣的問題，這才最終沒有成功。
實際上，朝鮮半島的情況也是比較類似的。但一方面，朝鮮語的音節結構畢竟比日語複雜，
而歷史政治原因又在其中作了重要推手，才使得朝鮮半島從混合系統走向了單一系統。
競爭如果最終導致走向單一系統，情況實際上就劃歸爲了另外兩種類型，因此也就不太需要過多討論了。
\section{正字法、文字改革與文字標準化}
在上面的討論中，雖然總涉及語文映射，只是抽象地討論映射與映射之間的差別，卻並沒有對具體的文字如何應用在具體的語言上進行探討。
而在這一章中，我們就要討論這一內容。一般來說，有關文字如何具體地應用到一門特定的語言上的內容，都可以認爲是\Term{zheng4zi4fa3}{正字法}{orthography}。
也就是說，正字法就是有關書寫系統的特徵的內容。研究正字法的學科可以稱作\Term{zheng4zi4fa3xue2}{正字法學}{orthology}

我們盡可能把面鋪開，將拼寫方法、書寫方向、標點用法等等，都歸結爲正字法。
本章也就對這些相關的內容進行總覽，然後討論有關文字改革和文字標準化的內容——雖然前者可以包括不同文字間的選取，但這兩者的主體應當都是在
同一門文字領域內，即正字法層面上的調整與規範，因此也放到這一章來。

\subsection{拼寫}
正字法最核心的內容是\Term{pin1xie3}{拼寫}{spelling}。這個詞對於表音文字來說更爲合適，對於表意文字來說，往往可能用「書寫」這個詞——
但方便起見，我們還是都統一成「拼寫」這個詞。我們首先從詞的角度來探討拼寫——當然，需要注意「詞」這個詞具有的模糊性。

\subsubsection*{表音文字}
理想狀況下，我們的文字應該是一種全音位文字，或者至少是元音附標文字。這樣，最理想的正字法當然是實現音位和字位的一一對應。
例如，現代俄語的спасибо，其音位和字位就是一一對應的。

但是，不會所有的情況都如此理想。
一種語言可能使用某些字位的組合來表示某個音位，這種情況多半出現在借詞或適應型文字中。
例如，拉丁語的ph、th、ch基本上都用來拼寫希臘語的借詞，對應希臘文的φ、θ、χ。當然，拉丁語原生就有一個字位組合，即qu，表示脣化的c。
在適應型文字中，使用字位組合的情況就更常見了，例如德語就中存在大量的sch、ch等字位組合。

當然，同一個字位或字位組合表示不同的音位也十分常見，但往往這種情況會通過語音環境或者拼寫環境區分。例如，德語的ch出現在a、o、u、au等元音後面時表示/x/，
其他情況下表示/ç/。不同的字位有時也用來表示同一個音位，這個情況就相對複雜了。這有可能是表示同一個音位的不同音位變體，但也有可能就是
表示歷史正字法的殘留。例如，梵語中所使用的\SA{न्}、\SA{ण्}、\SA{ङ्}，雖然其發音確實有差別，但實際上是同一個音位/n/的不同音位變體[n̪]、
[ɳ]、[ŋ]。而在現代希伯來語中，這種情況就更常見了，如\HE{ב}同時表示/b/和/v/，\HE{פ}同時表示/p/和/f/等等（不過人們有時
也會在字位中間加點區分）。 

還有一種情況是，一種語言的拼寫在某個時期定型了。隨着歷史的發展，語音不斷演變，拼寫卻幾乎沒有變化。例如，在古典藏語時期，藏文和藏語幾乎能夠
實現字位和音位的一一對應（當然，作爲元音附標文字，除了a），如\TI{སྒྲུབ}在古典藏語時期應當就是唸作sgrub——音位和字位一一對應。但是在現代
拉薩話中，這個詞讀作\IPA{/ʈ͡ʂup̚ $^{132}$/}——已經幾乎看不出字位與音位的對應關係了。

當然，最複雜的情況大概類似現代英語。英語除了歷史正字法之外，還有各種借詞直接套用了外語的正字法，導致很多詞已經不能從文字形式中看出發音了。
例如，daughter就是典型的歷史正字法殘留，古英語中這個詞寫作dohtor（現代德語中也有對應，即Tochter），h發音爲/x/。但在現代英語中，這個詞中間早就沒有了/x/的發音，
但文字形式中卻還殘留了一個gh。而cafe這個詞，按規則似乎應該讀作/keɪf/，但由於這是一個法語借詞，就得按來源讀作/ˌkæˈfeɪ/了。由於這種種
原因，英語幾乎沒有一套完全規則的正字法，而其發音必須一個個詞去記憶（當然，也不是完全沒有規律）。

\subsubsection*{表意文字}

表意文字的拼寫更加複雜，且幾乎每種表意文字都有自己獨特的特徵。

對於漢字來說，「拼寫」所反映的往往是字形上的區別，而不是表音文字那樣的字位組合方式上的區別。漢字本身經過長期的演化，其書體幾經迭代，
導致同字異形的情況相當普遍，其來源也相當複雜。

在漢字規範化以前，對同一個語素本身就存在部件多少、部件組合方式、飾筆添加與否以及書體風格
等多方面所導致的差異，如果這種差異沒有最後用來分別不同的字義，就會成爲同一個字位的全同異形字（也可以視爲字位變體），如「杯」和「盃」，
在任何情況下都不會有字義上的分別。另外，漢字從甲骨文、金文等古代書體變爲楷書這樣的現代書體，經過一個「隸定」的過程，而隸定方法的不同，
也會造成一批全同異形字，如「春」和「萅」的差別就純粹是隸定造成的。

除了這些不區分字義的情況，同一個語素對應的字可能在後世爲了區別字義而
通過增添部件等方法分化出幾個字出來，如「莫」一字分化爲「暮」和「莫」二字分擔了語義。而過去區分的字也有可能因爲某些原因合併爲一個字，
譬如在漢字簡化中所見到的「幹」、「乾」和「干」都簡化爲「\SC{干}」。

除了中文單一體系中的區別，漢字在不同地域的使用，也造成了字形的差別。
如對於同一個「漢」字，大陸、臺灣、日本、韓國字形分別爲「\SC{汉}」、「\TC{漢}」、「\JP{漢}」、「\KO{漢}」。

日本、朝鮮半島對漢字的運用，在某個時間段是與其本土文字（假名與諺文）混用的。黏着助詞在這些書寫系統中往往使用
本土文字，而漢字則用於漢語借詞的書寫，有時也用來書寫本土詞彙，即\Term{xun4du2}{訓讀}{kundoku}。

越南及中國境內部分少數民族對漢字的使用則與之不同，常常是用原生漢字書寫漢語借詞，而用漢字表音或者新造漢字來表示自己的原生詞彙。
如越南語的「\NOM{每𠊛調平等𠓀法律}」，其中「\NOM{每}」、「\NOM{平等}」、「\NOM{法律}」是漢語借詞，而「\NOM{𠊛}」、「\NOM{𠓀}」是原生詞，
「\NOM{調}」可能來自漢語。

部分漢字衍生的表意文字，如西夏文、女真文、契丹大字等，與漢字情況比較類似。但畢竟其歷史沒有那麼悠久，研究又沒有那麼充分，因此情況就相對好討論
一些。像是西夏文，就存在與漢字類似的異體字，如「\TG{𗀏}」和「\TG{𗍸}」。

而埃及聖書文和蘇美爾楔形文字等的情況與漢字相當不同，其並不是單純的表意文字系統，而是存在相當程度的混合成分。正如我們早已提到過的，這些表意文字
比起漢字，更像是混合系統。雖然我們可以在裏面找到非常多可以與漢字中的「字」對應起來的映射到態位的字位，可以稱爲\Term{yi4fu2}{意符}{ideographic component}，
但其中一部分字位可以單純僅僅充當表音的功能，我們這時稱其爲\Term{sheng1fu2}{聲符}{phonetic component}。
有一部分字位會跟聲符組合起來，限定其意義，我們稱其爲\Term{ding4fu2}{定符}{determinative}。
很容易看出，聲符和定符幾乎等同於漢字形聲字中的聲符和意符，但我們不會把漢字中的聲符和意符當作字位，
卻把這裏的聲符和定符會當作字位。原因在於，這些文字中的聲符和定符，似乎運用上比漢字自由得多
——當然，也許更多是出於習慣上的考慮。

蘇美爾語的楔形文字中，大多字位都能映射到態位，例如\CU{𒂵}ga表示「牛奶」、\CU{𒀭}an表示「天」、\CU{𒆕}du$_3$表示「建造」。
而這些字位又可以用作聲符，這裏面類似於音節文字，組成\CU{𒂵𒀭𒆕}表示gandu這個音，但是意義與「牛奶」、「天」卻沒有關係，而是表示「讓我建造」。
另外，\CU{𒆪𒃻}guza本身就可以表示「椅子」的意思，如果想要限定其含義，就可以在前面加上一個表示「木頭」的\CU{𒄑}ŋiš，組成\CU{𒄑𒆪𒃻}。
但\CU{𒄑}在\CU{𒄑𒆪𒃻}中並不發音，僅僅作爲定符限定意義。

埃及聖書文就更特殊了，我們可以在其中發現相當多個「部件」，而這些部件以或相鄰或堆疊的方式組合起來——我們習慣上就把這些部件當作字位，
雖然某種程度上來說會讓人有點膈應。聖書文中表示一個詞的，可以是單獨的意符，也可以是意符和聲符的組合，有時也會帶上定符。
聖書文中的聲符僅僅表示輔音，而且相對固定。有一部分聲符表示單個輔音，甚至可以整理出一套字母表，與輔音音位文字相當相似。
當然，聖書文也經常要在聲符旁加一個定符，如\begin{hieroglyph}{\leavevmode \loneSign{\Aca GO/35/}\HinterSignsSpace
\loneSign{\Aca GG/32/}\HinterSignsSpace
\loneSign{\Aca GD/92/}\HinterSignsSpace
\loneSign{\Aca GD/88/}}\end{hieroglyph}
表示「送」，其中前三個符號表音h\TIMES{Ꜣ}b，最後一個符號則是定符。
聖書文中的意符則常常通過加上一條豎線來表示，如\begin{hieroglyph}{\leavevmode \Cadrat{\CadratLineI{\Aca GO/32/}\CadratLine{\Aca GZ/32/}}}\end{hieroglyph}
，下面的豎線就表示上面的符號是意符，表示「房子」。
當然，聲符還有更複雜的用法。很多聖書文的聲符表示多個輔音，而又會加上表示單個輔音的聲符作爲補充（phonetic complement），最後還要加上定符。
如\begin{hieroglyph}{\leavevmode \Cadrat{\CadratLineI{\Aca GO/32/}\CadratLine{\Aca GD/52/}}\HinterSignsSpace
\loneSign{\Aca GD/88/}}\end{hieroglyph}
，\begin{hieroglyph}{\leavevmode \loneSign{\Aca GO/32/}}\end{hieroglyph}
本身就作爲聲符表示pr了，下面的\begin{hieroglyph}{\leavevmode \loneSign{\Aca GD/52/}}\end{hieroglyph}
r只是起到輔助提示發音的作用。
上面的敘述還是比較規則的情況，實際上的使用就更複雜了。

\subsection{斷詞、標點與方向}
除了拼寫之外，斷詞也是正字法相當重要的一部分。\Term{duan4ci2}{斷詞}{word seperating}指的是，在一種書寫系統中，如何通過\Term{kong4ge2}{空格}{space}等
方法來將有關單位區隔開來。

許多書寫系統都是需要斷詞的，例如我們所瞭解的英文、德文等等，但是其斷詞規則往往有所差異。大多書寫系統以詞爲基本單位用空格斷開，如英語
blue sky、a rainy day等等。但是，對於合成詞的處理，不同書寫系統存在差異。\Term{he2cheng2ci2}{合成詞}{compound}指的是不同詞幹組成的詞，
在英語中，許多合成詞都會把不同詞幹斷開，如high school等。但部分合成詞會不斷詞，如filename。也有一部分合成詞會使用\Term{lian2zi4fu2}{連字符}{hyphen}
「-」連接，如father-in-law。相比起來，德語的合成詞就更少斷詞，即便斷詞也會使用連字符，如Rindfleischetikettierungsüberwachungsaufgabenübertragungsgesetz
實際上可以斷成Rind-fleisch-etikettierungs-überwachungs-aufgaben-übertragungs-gesetz。

也有書寫系統並不按照「詞」來斷詞。如藏文會使用一個稱爲\TI{ཚེག}的符號「\TI{་}」將不同音節之間隔開，如\TI{བོད་ཡིག་}。而梵文的斷詞則相對更複雜，合成詞是
肯定不斷開的，但除此之外不同詞間可以通過
語音規則連接在一起，如\SA{नास्तीह}實際上是三個詞，即\SA{न अस्ति इह}。朝鮮語通常按照語法來斷詞，合成詞常常斷開，但黏着助詞一般不與實詞分開，
如「\KO{대한민국의 주권은 국민에게 있고}」，「\KO{대한민국}」、「\KO{주권}」、「\KO{국민}」等作爲實詞與後面的語法助詞相連。

另外，相當一部分書寫系統是不斷詞的。中文和日語都是典型的例子，但是偶爾也會使用「・」來斷開音譯詞。

\Term{biao1dian3fu2hao4}{標點符號}{punctuation}也是正字法的重要內容。最重要的一個標點符號就是\Term{ju4hao4}{句號}{full stop}，一般用來標識一句話的結束。
句號在不同書寫系統中有不同的形式，用法也相對有差異。譬如，大多數語言都是使用句號分開句子的，但中文中使用句號卻是按語義來判斷。
拉丁文、希臘文、西里爾文等等文字所對應的句號都是一個點「.」，而漢字（中文和日文）對應的句號都是一個圈「。」。相當多婆羅米系文字都會使用一條豎線「\SA{।}」
承擔類似句號的作用，泰文中更直接用空格分句，等等。至於其他的標點，各個語種的情況就相當不一致了。

\Term{fang1xiang4}{方向}{direction}在書寫系統中實際上有三個層面：一是\Term{zhou2}{軸}{axis}，指一行究竟是橫着的還是豎着的，
前者稱爲\Term{heng2shu1}{橫書}{horizontal writing}，而後者稱爲\Term{zong4shu1}{縱書}{vertical writing}；
二是\Term{hang2xu4}{行序}{lining}，指的是不同行的排列方向，基本上所有橫書的文字行序都是從上往下的，而縱書的文字則從左往右或從右往左都存在；
三是\Term{xiang4xing4}{向性}{directionality}，即一行中的文字是怎麼排列的，縱書的文字基本上都是從上往下排列的，而橫書的文字則有從右往左和
從左往右的區分。世界上絕大多數文字的方向都是橫書、行序從上往下、向性從左往右的，阿拉伯文、希伯來文等則通常是向性從右往左。
漢字系文字在過去基本都習慣縱書、行序從右往左，但現代也都會適應橫書。蒙古文等則是縱書、行序從左往右。
無論是行序還是向性，從下往上的文字是相當少的，僅有歐甘文（Ogham）等少數文字會違背這一規律。
埃及聖書文雖然方向總體上不定，但線性書寫與垂直堆疊並存，情況相對比較特殊。

在阿拉伯文、蒙古文等文字中，不同字位會通過特定的方式連接在一起，稱爲\Term{lian2xie3}{連寫}{joining}。
存在連寫的文字，其字位單獨出現與在詞中出現的樣貌相當不同。
如對於一個阿拉伯文字母的單獨形式\AR{ع}，其在詞中的連寫形式是\AR{ـعـ}，如在單詞\AR{العربية}中。
而蒙古文也類似，一個字母的單獨形式爲\MOV{ᠣ}，詞中連寫形式是\MOV{‍ᠣ‍}，如單詞 \MOV{ᠮᠣᠩᠭᠣᠯ}。

\subsection{拼音與轉寫}
不同的文字類型之間通過某種方式進行轉換，即\Term{zhuan3xie3}{轉寫}{transcription}。
轉寫在許多情景下具有相當重要的意義，在漢語語境下則一般指\Term{pin1yin1}{拼音}{pinyin}。

\subsection{文字改革與文字標準化}


\section{文字的排印與數位化}
